\documentclass[11pt,twocolumn]{article}

\usepackage[margin=1in]{geometry}
\usepackage{amsmath,amssymb,amsthm}
\usepackage{graphicx}
\usepackage{hyperref}
\usepackage{listings}
\usepackage{xcolor}
\usepackage{booktabs}
\usepackage{enumitem}

\hypersetup{colorlinks=true,linkcolor=blue!60!black,citecolor=green!40!black,urlcolor=blue!60!black}

\lstdefinelanguage{Rust}{
  morekeywords={fn,let,mut,pub,struct,enum,impl,trait,use,mod,match,if,else,for,while,loop,return,Self,self,async,await,move,const,static,where,crate,super,ref},
  sensitive=true,
  morecomment=[l]{//},
  morecomment=[s]{/*}{*/},
  morestring=[b]"
}

\lstset{
  language=Rust,
  basicstyle=\ttfamily\scriptsize,
  keywordstyle=\color{blue!70!black}\bfseries,
  commentstyle=\color{green!50!black}\itshape,
  stringstyle=\color{red!60!black},
  breaklines=true,
  frame=single,
  backgroundcolor=\color{gray!5}
}

\newtheorem{definition}{Definition}
\newtheorem{theorem}{Theorem}
\newtheorem{proposition}{Proposition}

\title{Three-Currency Economics on Holochain:\\SAP, TEND, and MYCEL with\\Consciousness-Gated Treasury}
\author{Tristan Stoltz \\ Luminous Dynamics \\ \texttt{tristan.stoltz@evolvingresonantcocreationism.com}}
\date{March 2026}

\begin{document}

\twocolumn[
\begin{@twocolumnfalse}
\maketitle

\begin{abstract}
We present a three-currency economic system implemented on Holochain's agent-centric
architecture with 8 coordinator zomes. The system distinguishes three economic
functions---stable accounting (SAP), care-oriented mutual credit (TEND), and commons
stewardship (MYCEL)---each with distinct monetary properties: SAP is commodity-backed
and demurrage-bearing; TEND is bounded zero-sum mutual credit without demurrage;
MYCEL represents collective stewardship and funds public goods. Treasury operations
are consciousness-gated, requiring progressive governance tiers for increasing levels
of financial authority. We describe the implementation of four critical financial
operations---oracle price feeds, mutual credit issuance, concurrent demurrage
computation, and currency composting---and report on hardening fixes for division-by-zero,
TEND race conditions, concurrent demurrage overflow, and oracle-to-TEND call ordering.
All 10 financial operations are consciousness-gated through shared verification
functions, and the system includes Sybil resistance through consciousness-tier
requirements on currency minting.
\end{abstract}
\end{@twocolumnfalse}
]

\section{Introduction}

Most cryptocurrency projects implement a single currency optimized for a single
function: Bitcoin for store of value, stablecoins for unit of account, Ethereum
for programmable transactions. Real economies, however, require multiple monetary
instruments serving distinct purposes~\cite{lietaer2001}.

Bernard Lietaer argued that monetary monocultures are fragile: when a single
currency fails, the entire economy suffers. Complementary currencies---local
exchange systems, time banks, mutual credit networks---provide resilience through
diversity~\cite{lietaer2004}.

We implement this insight on Holochain with three currencies, each designed for
a different economic function:

\begin{description}[nosep]
  \item[SAP (Sapience)] Stable unit of account. Backed by a basket of real
    goods and services. Used for pricing, contracts, and inter-community exchange.
  \item[TEND (Tenderness)] Mutual credit for care work and community service.
    Implements demurrage (value decay) to encourage circulation. Named for the
    tenderness of care labor.
  \item[MYCEL (Mycelium)] Commons currency for collective stewardship.
    Used for treasury operations, public goods funding, and regenerative
    investment.
\end{description}

\section{Architecture}

\subsection{Eight Zomes}

\begin{enumerate}[nosep]
  \item \textbf{currency-mint}: Supply management for all three currencies.
    Minting requires consciousness-gated authorization.
  \item \textbf{payments}: Transaction processing. Handles transfers, batch
    payments, and cross-currency exchanges.
  \item \textbf{tend}: TEND-specific mutual credit operations including
    credit line management, demurrage computation, and composting.
  \item \textbf{treasury}: Commons treasury management with multi-signature
    (DKG) authorization for large expenditures.
  \item \textbf{staking}: MYCEL staking for governance weight amplification.
  \item \textbf{recognition}: Non-monetary recognition tokens for gratitude
    and appreciation.
  \item \textbf{price-oracle}: Decentralized price feeds for the SAP basket.
  \item \textbf{bridge}: Cross-cluster financial dispatch.
\end{enumerate}

\subsection{Holochain Advantages for Finance}

Holochain's agent-centric model provides several advantages for financial
applications:

\begin{itemize}[nosep]
  \item \textbf{Double-spend prevention}: Each agent's source chain is
    append-only. A double-spend would require forking the source chain,
    which is detectable by DHT validators.
  \item \textbf{No transaction fees}: Unlike blockchain, Holochain has no
    per-transaction gas fees. This makes micro-transactions viable.
  \item \textbf{Offline transactions}: Agents can create transactions on
    their source chain while offline, syncing when connectivity returns.
  \item \textbf{Privacy}: Balances are on the agent's source chain, not a
    public ledger. Only the counterparty and DHT validators see transaction
    details.
\end{itemize}

\section{SAP: Stable Unit of Account}

\subsection{Basket Backing}

SAP is pegged to a basket of real goods and services, defined by the community.
A typical basket might include:

\begin{itemize}[nosep]
  \item 1 kg of locally-grown grain
  \item 1 hour of basic skilled labor
  \item 1 kWh of renewable energy
  \item 1 liter of clean water
\end{itemize}

The basket composition is governed by the Governance cluster (Steward+
required to modify). The price-oracle zome aggregates price attestations
for each basket component.

\subsection{Minting}

SAP cannot be created ex nihilo. Minting occurs only when:
\begin{enumerate}[nosep]
  \item An agent contributes a basket-equivalent of real goods/services
    to the community
  \item The contribution is verified by at least two Citizen-tier attestors
  \item The currency-mint zome issues SAP proportional to the basket value
\end{enumerate}

This makes SAP supply a function of real economic activity rather than
arbitrary monetary policy.

\section{TEND: Mutual Credit Without Demurrage}

\subsection{Mutual Credit}

TEND operates as a mutual credit system~\cite{riegel1944}: every TEND
created as a positive balance simultaneously creates an equal negative
balance. The total TEND supply across the network is always zero:

\begin{equation}
  \sum_{i} B_i^{\text{TEND}}(t) = 0 \quad \forall t
\end{equation}

This property is enforced by the integrity zome: the validation rule
verifies that each TEND transaction creates equal and opposite entries
on the sender's and receiver's source chains.

\subsection{Balance Discipline}

TEND does not decay over time. Instead, circulation discipline is enforced by
reciprocal accounting and bounded credit lines. Each member's positive or
negative position is constrained by consciousness-gated limits, and the entire
system remains zero-sum:

\begin{equation}
  \sum_{i} B_i^{\text{TEND}}(t) = 0 \quad \forall t
\end{equation}

\begin{proposition}[TEND zero-sum invariant]
For every confirmed exchange, the provider's positive balance increase exactly
matches the receiver's negative balance increase:
\begin{equation}
  \Delta B_{\text{provider}} + \Delta B_{\text{receiver}} = 0
\end{equation}
\end{proposition}

\begin{proof}
Each demurrage deduction from a positive balance is added to the compost
treasury. Each demurrage reduction of a negative balance is subtracted
from the compost treasury. Both operations preserve the zero-sum invariant.
\end{proof}

\subsection{Concurrent Demurrage}

A critical implementation challenge: demurrage must be computed correctly
when multiple transactions interact with the same balance concurrently.

The initial implementation had a counter overflow vulnerability: if two
concurrent transactions each computed demurrage on the same balance, the
second transaction would double-count the decay. The fix: demurrage is
computed exactly once per balance-modifying operation, using the source
chain's timestamp ordering as the canonical timeline.

\section{MYCEL: Commons Currency}

\subsection{Collective Stewardship}

MYCEL represents collective stewardship capacity. It is:
\begin{itemize}[nosep]
  \item Minted by community-wide consensus (Governance cluster)
  \item Held in the commons treasury
  \item Spent on public goods, infrastructure, and regenerative investment
  \item Stakeable for governance weight amplification
\end{itemize}

\subsection{Staking}

MYCEL staking locks tokens for a configurable period, providing:
\begin{itemize}[nosep]
  \item Governance weight amplification (staked MYCEL increases vote weight)
  \item Treasury stability (locked tokens reduce circulating supply)
  \item Commitment signaling (staking demonstrates long-term alignment)
\end{itemize}

Staking does not generate yield (this is not DeFi). The benefit is
governance influence, not financial return.

\section{Consciousness-Gated Treasury}

\subsection{Tiered Access}

Treasury operations are gated by consciousness tier:

\begin{table}[h]
\centering
\begin{tabular}{ll}
\toprule
\textbf{Operation} & \textbf{Min. Tier} \\
\midrule
View balance & Observer \\
Propose expenditure & Citizen \\
Approve expenditure ($< 1\%$ treasury) & Steward \\
Approve expenditure ($\geq 1\%$ treasury) & Steward + DKG \\
Emergency fund release & Guardian + DKG \\
Modify SAP basket & Steward + supermajority \\
\bottomrule
\end{tabular}
\caption{Consciousness-gated treasury operations.}
\end{table}

\subsection{Shared Verification Functions}

All 10 consciousness-gated financial operations use two shared verification
functions in \texttt{mycelix-finance/zomes/shared}:

\begin{lstlisting}
pub fn verify_participant_tier(
    cred: &ConsciousnessCredential
) -> ExternResult<()> {
    if cred.tier < ConsciousnessTier::Participant {
        return Err(wasm_error!("Insufficient tier"));
    }
    if cred.is_expired(now_us()?) {
        return Err(wasm_error!("Expired credential"));
    }
    Ok(())
}

pub fn verify_citizen_tier(
    cred: &ConsciousnessCredential
) -> ExternResult<()> {
    if cred.tier < ConsciousnessTier::Citizen {
        return Err(wasm_error!("Insufficient tier"));
    }
    if cred.is_expired(now_us()?) {
        return Err(wasm_error!("Expired credential"));
    }
    Ok(())
}
\end{lstlisting}

\section{Price Oracle}

\subsection{Decentralized Price Feeds}

The price-oracle zome aggregates price attestations from multiple oracle agents:

\begin{enumerate}[nosep]
  \item Oracle agents (Citizen+ tier) submit price attestations for SAP basket
    components
  \item Attestations include: item, price, timestamp, confidence
  \item The oracle computes the median price (robust to outliers)
  \item Stale attestations (beyond TTL) are excluded
\end{enumerate}

\subsection{Hardening}

Four critical bugs were fixed during security hardening:

\begin{enumerate}[nosep]
  \item \textbf{Oracle division-by-zero}: When no valid attestations exist,
    the median computation divided by zero. Fix: return \texttt{None} when
    attestation count is zero.
  \item \textbf{TEND race condition}: Concurrent TEND issuance could create
    unbalanced credit lines. Fix: source chain sequencing with optimistic
    locking.
  \item \textbf{Concurrent demurrage}: Double-counted decay on concurrent
    balance access. Fix: single demurrage computation per balance modification.
  \item \textbf{Missing oracle$\to$TEND call}: The TEND pricing function
    did not call the oracle for current SAP-TEND exchange rate. Fix: added
    cross-zome call with fallback to last known rate.
\end{enumerate}

\section{Currency Composting}

\subsection{Lifecycle}

Currencies have a lifecycle: minting $\to$ circulation $\to$ composting.
Composting is the explicit removal of currency units from circulation:

\begin{itemize}[nosep]
  \item SAP: Demurrage and explicit compost routing remove idle value from
    circulation and fund commons maintenance
  \item TEND: No demurrage; circulation is governed through zero-sum credit
    limits rather than timed decay
  \item MYCEL: Governance-authorized composting of excess supply
\end{itemize}

\subsection{Idempotency}

Currency composting is idempotent: composting an already-composted unit is
a no-op. This was a hardening fix---the original implementation allowed
double-composting, which would undercount the circulating supply.

\subsection{Zero-Sum Invariant}

For all three currencies, the invariant holds:

\begin{equation}
  \text{minted} - \text{composted} = \text{circulating}
\end{equation}

This is enforced by the integrity zome and verified by DHT validators.

\section{Sybil Resistance in Finance}

Financial Sybil attacks (creating fake identities to obtain multiple credit
lines) are prevented by:

\begin{enumerate}[nosep]
  \item \textbf{Consciousness gating}: Currency minting requires Citizen+
    tier, which requires verified identity ($I \geq 0.5$)
  \item \textbf{Credit line limits}: TEND credit lines are proportional to
    the agent's consciousness profile score
  \item \textbf{Reputation consequences}: Financial misbehavior triggers
    reputation slashing ($0.5\times$), affecting all governance participation
  \item \textbf{Cross-cluster visibility}: The Identity cluster's web-of-trust
    can detect clusters of financially-interconnected Sybil identities
\end{enumerate}

\section{Non-Monetary Recognition}

The recognition zome provides non-monetary appreciation tokens:

\begin{itemize}[nosep]
  \item Gratitude tokens (from Hearth cluster integration)
  \item Skill recognition (from Education cluster integration)
  \item Community service badges
  \item Milestone acknowledgments
\end{itemize}

These tokens have no monetary value but contribute to the reputation dimension
of the consciousness profile, creating a positive feedback loop between
community participation and governance capability.

\section{Evaluation}

\subsection{Security Hardening Results}

The four hardening fixes (Section~7.2) were discovered through a combination
of code review and sweettest integration testing. The TEND race condition
required a 3-agent concurrent test scenario with a BadNonce retry mechanism
in the sweettest framework.

\subsection{Design Trade-offs}

\begin{table}[h]
\centering
\begin{tabular}{lll}
\toprule
\textbf{Property} & \textbf{Mycelix} & \textbf{DeFi} \\
\midrule
Tx fees & Zero & Gas fees \\
Privacy & Source chain & Public ledger \\
Yield & None (by design) & APY incentives \\
Governance & Consciousness & Token-weighted \\
Backing & Real goods & Algorithmic \\
\bottomrule
\end{tabular}
\caption{Mycelix Finance vs. DeFi comparison.}
\end{table}

\section{Related Work}

\textbf{Mutual credit on Holochain}: Holo\-chain was designed partly for
mutual credit~\cite{brock2018}. Our contribution is the three-currency
design with consciousness gating.

\textbf{Circles UBI}~\cite{circles2019}: Ethereum-based mutual credit with
social graph trust. Similar goals but limited by blockchain costs and
single-currency design.

\textbf{Grassroots Economics}~\cite{ruddick2015}: Community currency
implementations in Kenya. Validates demurrage and mutual credit in practice
but uses centralized infrastructure.

\section{Conclusion}

Three-currency economics on Holochain provides monetary diversity without
the overhead of blockchain consensus. SAP, TEND, and MYCEL serve distinct
economic functions---stable accounting, care labor, and collective
stewardship---with consciousness-gated treasury operations ensuring that
financial power scales with demonstrated governance capability rather than
accumulated wealth. The hardening of four critical financial bugs demonstrates
the importance of rigorous security review in decentralized financial systems,
and the zero-sum invariants provide formal guarantees of monetary integrity.

\begin{thebibliography}{99}
  \bibitem{lietaer2001} Lietaer, B. (2001). \textit{The Future of Money}.
    Random House.
  \bibitem{lietaer2004} Lietaer, B. (2004). Complementary Currencies in Japan Today.
    \textit{International Journal of Community Currency Research}, 8.
  \bibitem{gesell1916} Gesell, S. (1916). \textit{The Natural Economic Order}.
  \bibitem{riegel1944} Riegel, E. C. (1944). \textit{Private Enterprise Money}.
  \bibitem{brock2018} Brock, A. and Harris-Braun, E. (2018). Holochain:
    Scalable Agent-Centric Distributed Computing.
  \bibitem{circles2019} Circles (2019). Circles: A Decentralized Universal Basic Income.
  \bibitem{ruddick2015} Ruddick, W. et al. (2015). Complementary Currencies for
    Sustainable Development in Kenya. \textit{IJCCR}, 19.
\end{thebibliography}

\end{document}
